\chapter[Music's Meaning beyond Concepts]{How Can Music Mean without Images or Concepts?}
\section{An Earphone That Does Not Sound}
First, pick up an object: an old pair of earphones.

The left driver has broken; only the right works. During a ten-minute break, you put it in and play the song everyone knows. Familiar melody fills one ear; the other receives chairs scraping, homework shouted, cicadas outside. Strangely, even half the sound changes your heartbeat, loosens your shoulders, and summons images never sung: last month's sports day, perhaps, or an afternoon never returning.

Remove it and examine this bundle of waves. It has no images: no houses, faces, landscapes drawn. No concepts: it says neither friendship, parting, nor justice; no word appears. A dictionary cannot define it, a painting's label cannot attach to it. Something without pictures or words should seemingly pass like cicadas' noise, heard and gone.

Yet it can bring tears close.

Painting speaks through images, literature through words. Music seems empty-handed yet may reach the heart most directly and least reasonably. Where does meaning live without images or concepts? Did the composer put it there, the listener bring it, or are we sentimental animals projecting onto waves?

We will take those invisible waves apart.

\section{Four Answers in Music Class}
Enter a scene.

First afternoon lesson, late spring, slow ceiling fans in an ordinary middle-school music room. Half-drawn curtains leave bright bands on the floor. Ms. Chen, teaching music for twenty years, leaves the textbook closed and approaches an old speaker beside the piano. "No singing yet. Listen to a piece without words and tell me what you hear."

A cello sounds: slow, low, long-tailed notes, pauses slightly longer than expected. The room quiets; even the boy spinning his pen at the back stops.

Two minutes later, silence. "What did you hear?"

The literature class representative raises her hand. "Sadness. Someone saying goodbye, unwilling but needing to leave." She sounds certain, as though retelling a text.

The physics enthusiast who dismantles everything frowns. "I didn't 'hear' anything. Frequencies vibrating, fast or slow, high or low. Calling it sad is like calling a thermometer's mercury gloomy."

A quiet girl speaks softly. "Not sadness... I thought of my grandmother. She died last year. I don't know why this piece reminded me; it has nothing to do with her." She sounds puzzled herself.

The boy at the back gives the fourth answer. "It tricked me. Several times I expected a note to land, and it circled away before returning. It pulled me along. I liked it."

Four answers, one cello. Ms. Chen writes sadness, frequencies, grandmother, tricked, then asks our real question:

"Which of these did the music itself say?"

Remember this classroom. The chapter answers that apparently ungradable question.

\section[How Can Sound Produce Feeling?]{Where Does Feeling Come from without Images?}
Lay out the conflict before answering.

The physics enthusiast sounds most scientific. Sound is air vibrating; frequencies are measurable. Cello and electric drill differ physically in parameters, not meaning. An analyzer finds no component called sadness. Emotional expression seems projection onto nature, like melancholy mercury.

The literature representative has a firm case too. Why do slow, low, descending tunes sound heavy to almost everyone, quick, high, leaping tunes light? Nobody worldwide hears a funeral march as celebration. Such agreement cannot be billions independently projecting. Composers write through tears or smiles, performers sustain breath and pressure; did all these emotions vanish until listeners invented them?

The grandmother answer complicates matters. Meaning may lie between music and your memories, not composer and music. An unrelated cello opens her personal drawer. Is this music's meaning, or is music a master key accidentally opening different locks?

The fourth answer points elsewhere. Perhaps music manipulates rather than expresses, making expectations, delaying them, fulfilling or cheating them. Enjoyment resembles a story holding you in suspense: an ingenious machine, not a confiding heart.

Behind the positions lies a real question, not a contest in artistic sensitivity:

\textbf{What kind of meaning can waves without images or concepts have for people?}

Cargo in sound for every hearer? A composer's emotional parcel delivered to audiences? Listeners' private contraband? Automatic responses to human mechanisms? Choose provisionally which classroom answer best deserves the phrase "the music's meaning."

\section[Expression or Induced Emotion?]{Does It Express Emotion or Give You Emotion?}
Strengthen both positions before examining other cultures. Calling one side unmusical and the other sentimental is unfair; give your opponent reasons fully.

\textbf{First: music is the sound of emotion.}

Three reasons support the literature representative. Emotion already connects to voice: crying slows and lowers it, laughter quickens and raises it, anger makes it loud and forceful. Music refines, extends, and weaves this bodily language rather than merely symbolizing it. Second, untrained children largely agree which of two pieces sounds happy or sad, difficult for pure projection to explain. Third, composition and performance are embodied: breath, pause, pressure in a sustained note speak bodily. Denying emotion resembles calling a choked-up story merely sound waves.

\textbf{Second: you bring emotion; music supplies sound.}

Give the physicist three reasons too. Responses differ even to supposedly sad music: sorrow, grandmother, structural pleasure. Why does emotional cargo fail to arrive identically? Cross-cultural evidence is sharper. Western minor modes, roughly darker scale organizations, suggest sadness and major brightness, but centuries of churches, courts, and concerts trained that association. Other systems loosen or defeat it. Finally, the same piece means happiness at a wedding, summons on a battlefield, sales in an advertisement. Context-following emotion cannot live intrinsically in notes.

Consider New Orleans jazz funerals, a tradition over a century old. Bands play slowly and heavily toward burial, then brightly on return, with dancing along the road. Outsiders may mistake disrespect or frivolity. Within the tradition, release of the soul allows the living to continue: joy is an essential half of mourning, not its absence. Cultures organize musical emotion differently. Your intuition is habit, not a ruler for measuring others' feeling.

How are such habits made? Familiar pop intervals rely on a mathematical convention, twelve-tone equal temperament, dividing the octave equally into twelve. Zhu Zaiyu calculated its values with extraordinary precision in sixteenth-century Ming China using an abacus; Europe followed decades later, and Bach's Well-Tempered Clavier became a masterpiece within this grid. But not everyone uses it. Javanese and Balinese gamelan, ensembles centered on metal percussion, use differently spaced scales. Western listeners may hear wrong tuning; gamelan musicians may find piano awkward. Neither ear is wrong: each grew in different sonic water. Listening resembles accent, mistaking hometown speech for universal standard.

Each position holds half the truth. Music borrows human emotional embodiment, but culture largely teaches particular sound-emotion mappings. Can we inspect its components without mystical musical instinct?

\section{A Bridge for Thought: Five Layers of Music}
Yes. Any culture's music can be examined through five components, not examination facts but portable tools. Try them on pop, game scores, the call to prayer, or public dance speakers.

\textbf{First: pitch, the step a sound occupies.}

Pitch rises with vibration frequency. Imagine stairs, each note a step; a tune traces a path, climbing gradually, leaping five steps, lingering. Ms. Chen's cello feels heavy partly because it stays low and often descends.

\textbf{Second: rhythm, the shape of time.}

Duration and density form music's skeleton. Your heartbeat and footsteps are two percussion instruments. Quick, crowded rhythms recall running and excitement; sparse, slow ones walking and sighs. Accent matters too: "BOOM-da-da" and "da-da-BOOM" use the same three sounds but suggest heartbeat or waltz. Dances from Xinjiang's meshrep to Argentine tango supply bodily language for particular accent patterns.

\textbf{Third: harmony, colors layered together.}

Several simultaneous notes make harmony. One note is a color; combinations feel clear and settled, like blue and white, or tense and abrasive, like clashing fluorescents. Tension creates a desire for resolution, ears awaiting stability. This journey powers much music.

\textbf{Fourth: timbre, sound's face.}

One pitch and duration played on flute, tapped on a pencil case, or hummed by a voice has three faces. Timbre is rough, clear, ragged, smooth character. You recognize a singer in a second through timbre, not pitch. Traditions explore it through electric-guitar distortion, guzheng glissandos, low Tibetan-operatic throat tones.

\textbf{Fifth: structure, repetition and change.}

Music builds through time: what comes first, returns, changes in returning. Repetition lets you recognize routes; variation prevents boredom. Verse--chorus cycles, folk call-and-response, classical themes broken, inverted, accelerated, reassembled share this machinery. The boy's trick was structure establishing and breaking rules.

Return with these tools. Physics measures first- and second-layer parameters. Farewell emerges from low pitch, slow rhythm, descent, a tired body's acoustic outline. Trickery combines harmonic and structural expectations. Grandmother remains beyond all five: the tool explains what sound does, not why her particular drawer opened. Remember this boundary.

\section[An Ancient Report on Music]{A Music Report from Two Thousand Years Ago}
Read primary material. Ancient China already formally investigated sound's meaning.

The Record of Music in the Book of Rites, developed from the Warring States into early Han, opens directly:

\begin{quote}
All musical sound arises from the human heart. The heart is moved by external things. Stirred through contact with things, it takes form in sound.
\end{quote}
Music arises inwardly because the world moves the heart, whose movement becomes audible.

Notice the chain: external things $\rightarrow$ heart $\rightarrow$ sound. Music falls neither from heaven nor grows from instrument wood; a world-touched heart flows outward in sound. The text distinguishes sheng, individual sound; yin, sounds ordered by pattern; and yue, patterned sounds with dance in ritual. Music comes bundled with movement, ceremony, and social order, not merely sound art.

The Han-era Great Preface to the Book of Songs extends the chain into politics:

\begin{quote}
The music of an ordered age is peaceful and joyful, its government harmonious; that of disorder resentful and angry, its government awry; that of a ruined state mournful and troubled, its people distressed.
\end{quote}
Peaceful politics produces calm joy, disorder anger, collapse sorrow. Before judging truth, inspect its operation: music becomes society's dashboard, an era heard through its songs. Elders deploring today's tunes use this same two-thousand-year-old instrument.

The Analects, "Transmission," offers livelier evidence:

\begin{quote}
Hearing Shao music in Qi, the Master did not taste meat for three months. He said, "I never imagined music could reach such heights."
\end{quote}
Shao, attributed to legendary Shun, so absorbed Confucius that meat lost flavor for months. One of the earliest musical reviews admits sound can displace other senses. He analyzes neither scales nor harmony but reports absorption, akin to your near-tears during ten minutes with earphones.

These witnesses need not be invariably right. They establish our classroom dispute's ancient record. Since the Record of Music, Chinese thought linked music to hearts, rituals, and social order rather than pure sound. Whether that bundle helps or harms remains for comparison.

\section{One Melody, Three Explanations}
Return to the cello. Separate evidence, two minutes of slow low descending music and four answers; causes, our debate; students' own understanding, their reports; and judgment, who heard correctly. Compare causes.

\textbf{First: music borrows the body's acoustic contours.}

Slow, low, descending sound resembles tiredness, weakness, crying; quick, high, dense sound resembles excitement, running, laughter. Bodies synchronize: breathing slows, shoulders sink, tears prepare. This straightforward account explains shared intuitive responses and composers describing tearful creation. It handles embodied sadness and joy, but not the boy's pleasure in structure, which imitates no particular bodily condition.

\textbf{Second: music is an expectation machine.}

The grip lies in prediction, not sound alone. Three rising notes invite a fourth, tension promises resolution. Musical craft plays with bets: immediate payoff pleases, delay grips, deception surprises, deception followed by fulfillment pleases most. This explains structural fascination, boredom after predictable songs, and incomprehension when complexity prevents prediction. But grandmother's memory is not puzzle-solving; music seems to bypass prediction and enter recollection directly.

\textbf{Third: occasions and memories attach meaning.}

This sociological account sees meaning repeatedly attached through settings, companions, rituals until sounds carry their scent and learning feels natural.

Count your own attachments. A birthday tune instantly cues action. Originally a late-nineteenth-century American kindergarten greeting, global parties made it a ritual switch. A wedding march originally from Wagner's Lohengrin became a signal for tears through repeated ceremonies. A national anthem's pitches and rhythms acquire a community's weight through flags, victories, ceremonies, and required standing. Without that process, an anthem has no acoustic difference from another march.

This also explains music's ancient work in ritual, dance, identity. Gamelan drives Javanese shadow theater and court ceremony, not merely accompanies it. In many West African traditions, drums speak names and summon villages; dancers and drummers converse rhythmically. Dong grand song in Guizhou is unaccompanied multipart singing without a conductor; singing groups organize village social identity through membership and exchanges. School sports entrances use marches to coordinate feet, not just please ears. A generation's shared song likewise carries shared occasions.

The limit: first hearings without associations can still give goosebumps. Attachment explains a birthday song's tears, not why this melody rather than another was selected for global attachment.

Each account holds part: bodily contour intuition, expectation structure, social attachment culture and memory. Use all three rulers instead of choosing a camp.

\section[Is Music's Content Only Sound?]{A Deeper Challenge: Is Music's Content Only Sound?}
Now introduce Eduard Hanslick, nineteenth-century Austrian critic whose slim On the Musically Beautiful of 1854 began a continuing argument.

He challenged the assumption that music expresses emotion as literature tells stories. Definite emotions, he argued, have objects: sadness over something, joy for someone. Remove the object and only physiological excitement remains. Music supplies no such object; a melody cannot state whom or why it mourns. It can arouse listeners but not articulate a determinate feeling. Its own content is moving tonal forms, notes pursuing, overlapping, developing like a kaleidoscope's patterns, about nothing else yet beautiful and ordered. Formalism makes organization itself content without importing emotions.

State his reasons fairly; he is not a cold person forbidding feeling. He denies identifying arousal with content, not arousal itself. Chilies produce sweat and tears, but these are responses, not ingredients; capsaicin is an ingredient. Likewise sadness belongs to your response, not notation. This frees composers from pretending to speak all humanity's emotions and listeners from feeling incompetent if no sadness appears. They can simply love sound's movement. The same allegedly sad piece gives others beauty or grandeur while the score remains silent.

Correct an easy misreading: formalism does not prove music meaningless or understanding impossible. It relocates meaning from emotional translation to sound's organization. Fugue voices pursuing one another, inversion, extension, fragmentation, reassembly are real, identifiable, debatable events belonging to music, not private fantasy. Understanding follows what sound does rather than translating a mood. On this definition, the boy hearing trickery understood best.

Its limits remain. Pure instrumental music fits well, but songs with words, funerals, and public ceremonies exceed form alone. Birthday music works partly because ordinary form lets everyone sing. Grandmother's memory likewise happens between listener and melody, neither entirely score nor mind. Formalism excellently maps what sound does, but occasion, memory, identity, and ritual lie beyond its coast. A theory earns value by drawing identifiable boundaries, not saying everything.

\section{The Playlist Age: Who Shapes Your Ears?}
This ancient question holds court daily in your phone.

Inspect unquestioned app categories: classical, pop, folk, rock, Chinese-style, electronic, arranged naturally as supermarket shelves. Remove the shelving. These three or more categories were cut by historical institutions, not music itself.

Classical music is institutionally made. Bach was a salaried church employee delivering weekly worship work. Nineteenth-century concert institutions placed his pieces before formally dressed, silent audiences. Conservatories and graded examinations continued, turning refinement into a measurable, enrollable, fee-paying qualification from grade one to ten. A child's piece often counts as classical because it appears in examination books rather than through sound alone.

Folk music was similarly defined by its opposite. Field songs existed unnamed for centuries until academic collectors arrived with notebooks and recorders, transcribed, archived, and staged them. Recording also transformed: multipart textures simplified, dialect corrected, improvisation fixed into standard versions. Textbook folk songs often come as polished production-line editions.

Technology bounded pop. Early-twentieth-century records held three or four minutes per side, compressing songs into that span, still largely familiar today. Not emotion's natural duration but a century-old wax disc's capacity. Broadcasting and charts then determined what could play and rank as popular.

One example crosses every boundary. You may sing "Beyond the farewell pavilion, beside the ancient road, fragrant green grass meets the sky"---Farewell. Its melody came from American composer Ordway, traveled to Japan for new lyrics, then Li Shutong supplied Chinese words for modern school singing. Chinese or American, classical, folk, or pop? No box holds it neatly. Its birth certificate lists institutions: schools, printed scores, international circulation. History draws boundaries; melody does not grow them.

Algorithms now hold the pen. A platform knows neither Beethoven nor Dong grand song, only when you skip. Short videos attach fifteen-second choruses to millions of clips, making a whole song's fate depend on them. Creators front-load hooks for impatient algorithmic ears. Not wholly bad: small-town dialect rappers can bypass stations and labels to reach national listeners. But recognize how much of supposed personal taste follows a recommendation system's feeding formula.

Ritual attachment continues nearby. School bells and exercise music tie discipline to melody; New Year mall songs tie shopping to wishes; game and stream battle themes tie adrenaline to beat. They do the Record of Music's old work, organizing hearts, with property managers, platforms, and operating teams replacing priests and courts. You need not be angry, but should notice who pressed play, when, and why.

\section{Your Question: If a Music Disappeared}
Return to the classroom's sadness, frequencies, grandmother, tricked. Each holds truth, none all of it.

End with a larger question without a standard answer.

Dong grand-song masters in Guizhou villages are aging. Unconducted, unwritten, orally transmitted multipart singing requires dozens growing up together. Fewer youths learn because its life, growing, working, and exchanging songs together in drum towers, is disappearing, not because it sounds bad. Suppose the last group dissolves while recordings and transcriptions remain playable in museums.

What has humanity lost?

One answer: nothing. Waves and spectra remain on drives, reproducible and analyzable through five layers. Another: everything. The song was groups, towers, festivals, dialect, village memory joined together, not stored waves. Sound can be archived, its life cannot. A middle answer: some lost, some saved, depending on which layer you mean.

Give reasons using all our tools. Which layers survive recording? Which bodily and predictive meanings might outlive occasions? Which social attachments disappear with life? If Hanslick's form-as-content is right, perhaps the core remains. If the Record of Music's world--heart--sound chain is right, can sound remain music after the heart is gone?

Your own ten-minute listening sessions participate too. Today's repeated song attaches memories for you ten years hence. Without images or concepts, music may be humanity's best container for "we once lived this way." What it holds and who decides: you now have ears able to take it apart.
